% diagrams/firstwords/abuela-puerta.tex -- ¡en el pueblo! Grandma espera
% en la puerta de su casa; en el patio, el limonero y el gallinero.
\begin{tikzpicture}[dibujofw]
  \begin{scope}[scale=1.3]
    % el gallinero, con su rampa
    \filldraw[fill=white] (-7.8,0.6) rectangle (-5.4,2.6);
    \filldraw[fill=white] (-8.2,2.5) -- (-6.6,3.6) -- (-5.0,2.5) -- cycle;
    \filldraw[fill=white] (-7.0,1.0) rectangle (-6.2,2.0);
    \filldraw[fill=white] (-7.0,1.0) -- (-6.2,1.0) -- (-5.3,0) -- (-6.1,0) -- cycle;
    \foreach \x in {-7.7,-5.8} {\filldraw[fill=white] (\x,0) rectangle ({\x+0.3},0.6);}
    % la casa: la pared, el tejado, la puerta abierta y la ventana con flores
    \filldraw[fill=white] (-4.2,0) rectangle (3.8,6.8);
    \filldraw[fill=white] (-4.8,6.6) -- (-0.2,9.2) -- (4.4,6.6) -- cycle;
    \foreach \y in {7.3,8.0} {\draw ({-4.8+4.6*(\y-6.6)/2.6},\y) -- ({4.4-4.6*(\y-6.6)/2.6},\y);}
    \filldraw[fill=white] (-1.4,0) rectangle (1.4,5.3);
    \filldraw[fill=white] (1.4,0) -- (2.3,0.35) -- (2.3,4.95) -- (1.4,5.3) -- cycle;
    \filldraw[fill=white] (-3.6,3.0) rectangle (-2.0,4.9);
    \draw (-2.8,3.0) -- (-2.8,4.9);
    \filldraw[fill=white] (-3.8,2.6) rectangle (-1.8,3.0);
    \foreach \x in {-3.4,-2.8,-2.2} {\foreach \a in {0,90,180,270} {\filldraw[fill=white] ([shift={(\a:0.16)}]\x,3.2) circle (0.12);}}
    % el limonero
    \begin{scope}[shift={(6.4,0)}, scale=0.85] \limoneroFW \end{scope}
    \draw (-8.8,0) -- (9.2,0);
  \end{scope}
  % Grandma, en la puerta, saludando
  \begin{scope}[scale=0.8] \mujerFW{Grandma}{Abajo}{Arriba} \end{scope}
\end{tikzpicture}
